\documentclass[11pt,a4paper]{article}

% Packages
\usepackage[utf8]{inputenc}
\usepackage[T1]{fontenc}
\usepackage{amsmath, amssymb, amsfonts}
\usepackage{geometry}
\geometry{a4paper, margin=2.5cm}
\usepackage{xcolor}
\usepackage{tcolorbox}
% Rimosso 'shadows' per evitare errori di compilazione
\tcbuselibrary{skins, breakable, theorems}
\usepackage{hyperref}
\usepackage{graphicx}
\usepackage{enumitem}

% Colors setup
\definecolor{sapienzared}{RGB}{130, 36, 51}
\definecolor{theoryblue}{RGB}{20, 60, 120}
\definecolor{practicegreen}{RGB}{30, 100, 50}

% Custom Box Environments
\newtcolorbox{theorybox}[1]{
    colback=theoryblue!5,
    colframe=theoryblue,
    fonttitle=\bfseries,
    title=#1,
    arc=3mm,
    breakable
}

\newtcolorbox{examtip}{
    colback=sapienzared!5,
    colframe=sapienzared,
    fonttitle=\bfseries,
    title=Exam Tip / Watch out!,
    arc=1mm,
    leftrule=3mm
}

% Title
\title{\textbf{\Huge Robotics 2: The Ultimate Midterm Guide} \\ \Large Theory \& Step-by-Step Exercise Solver}
\author{Based on Prof. De Luca's Lectures}
\date{\today}

\begin{document}

\maketitle
\tableofcontents
\newpage

% ==========================================
% PART 1: THEORY
% ==========================================
\part*{Part 1: Theory Core Concepts}

\section{Kinematic Calibration}
Calibration improves robot accuracy by updating nominal Denavit-Hartenberg (DH) parameters using external measurements. The actual end-effector position $p$ depends on the parameters $\xi = [\alpha, a, d, \theta]^T$.
Linearizing the direct kinematics around the nominal parameters yields the basic calibration equation:
\begin{equation}
    \Delta p = \Phi(q) \Delta \xi
\end{equation}
where $\Delta p = p_{measured} - p_{nominal}$ is the Cartesian error, $\Delta \xi$ is the vector of parameter errors, and $\Phi(q)$ is the \textbf{Regressor Matrix}. 

\begin{theorybox}{Building the Regressor Matrix}
The matrix $\Phi$ is built by taking the partial derivatives of the direct kinematics $k(q)$ with respect to the specific uncertain parameters, evaluated at their nominal values:
\begin{equation}
    \Phi_i(q) = \frac{\partial k(q)}{\partial \xi_i}\Bigg|_{nom}
\end{equation}
For multiple measurements at different configurations $q_1, q_2, \dots, q_N$, we stack the equations and solve using the pseudoinverse: $\Delta \xi = \Phi_{stacked}^{\#} \Delta p_{stacked}$.
\end{theorybox}


\section{Kinematic Redundancy}
A robot is redundant if the dimension of the joint space $n$ is strictly greater than the dimension of the task space $m$ ($n > m$).
The differential kinematics is $J(q)\dot{q} = \dot{x}$. Since $J$ is $m \times n$, there are infinite solutions. 

\begin{theorybox}{Redundancy Resolution Methods}
\begin{itemize}
    \item \textbf{Pseudoinverse (Minimum Norm):} $\dot{q}_{PS} = J^\# \dot{x} = J^T(JJ^T)^{-1}\dot{x}$. Minimizes $\|\dot{q}\|^2$.
    \item \textbf{Weighted Pseudoinverse:} $\dot{q}_W = W^{-1}J^T(J W^{-1} J^T)^{-1}\dot{x}$. Minimizes $\dot{q}^T W \dot{q}$.
    \item \textbf{Null-Space Projection:} $\dot{q} = J^\# \dot{x} + (I - J^\# J)\dot{q}_0$. The term $(I - J^\# J)$ projects any desired velocity $\dot{q}_0$ into the null space of $J$ (self-motion).
    \item \textbf{Inertia-Weighted Pseudoinverse (Minimizes Kinetic Energy):} If we choose the weight matrix $W = M(q)$ (the inertia matrix), the weighted pseudoinverse $\dot{q}_M = M^{-1}J^T(J M^{-1} J^T)^{-1}\dot{x}$ minimizes the instantaneous kinetic energy $T = \frac{1}{2}\dot{q}^T M(q) \dot{q}$.
\end{itemize}
\end{theorybox}


\section{Lagrangian Dynamics}
The dynamic model of a rigid manipulator is derived from the Lagrangian $\mathcal{L}(q, \dot{q}) = T(q, \dot{q}) - U(q)$, where $T$ is kinetic energy and $U$ is potential energy. The Euler-Lagrange equations yield:
\begin{equation}
    M(q)\ddot{q} + c(q, \dot{q}) + g(q) = \tau
\end{equation}
\begin{itemize}
    \item $M(q)$: $n \times n$ symmetric, positive definite Inertia Matrix.
    \item $c(q, \dot{q})$: Coriolis and centrifugal forces vector. $c(q,\dot{q}) = S(q,\dot{q})\dot{q}$.
    \item $g(q)$: Gravity vector. $g(q) = (\frac{\partial U}{\partial q})^T$.
\end{itemize}

\subsection{Actuator Dynamics and Friction}
In a real manipulator, torques are provided by motors through reduction gears. If $n_{ri}$ is the gear ratio of the $i$-th joint and $I_{mi}$ is the motor rotor inertia, the kinetic energy of the motors adds a constant diagonal matrix to the inertia matrix:
\begin{equation}
    M_{tot}(q) = M(q) + \text{diag}(n_{ri}^2 I_{mi})
\end{equation}
Additionally, the real torque $\tau$ must compensate for joint friction. The complete model becomes:
\begin{equation}
    M_{tot}(q)\ddot{q} + c(q, \dot{q}) + g(q) + F_v \dot{q} + F_c \text{sgn}(\dot{q}) = \tau
\end{equation}
where $F_v$ is the viscous friction diagonal matrix and $F_c$ is the Coulomb friction diagonal matrix.

\subsection{Important Properties of the Dynamic Model}
\begin{enumerate}
    \item \textbf{Skew-Symmetry:} The matrix $\dot{M}(q) - 2S(q, \dot{q})$ is skew-symmetric IF $S$ is derived using Christoffel symbols. This implies $x^T(\dot{M}-2S)x = 0 \quad \forall x$.
    \item \textbf{Linear Parameterization:} The dynamic model is linear with respect to a set of constant dynamic parameters (masses, lengths, inertias). It can be written as $Y(q, \dot{q}, \ddot{q})\mathbf{a} = \tau$, where $Y$ is the dynamic regressor and $\mathbf{a}$ is the vector of dynamic coefficients.
\end{enumerate}

\newpage
% ==========================================
% PART 2: HOW TO SOLVE EXERCISES
% ==========================================
\part*{Part 2: How to Solve Midterm Exercises}
\textit{This section provides step-by-step algorithms to solve the exact typologies of exercises found in the exams.}

% ---------------------------------------------------------
\section{Type A.1: Task Priority and Task Augmentation}

{Algorithm: Solving Task Priority (TP) Problems}
\textbf{Goal:} Execute a primary task $\dot{x}_1$ perfectly, and a secondary task $\dot{x}_2$ as best as possible using the remaining degrees of freedom.
\begin{enumerate}
    \item Identify the Jacobians $J_1$ and $J_2$.
    \item Calculate the pseudoinverse of the primary task: $J_1^\# = J_1^T(J_1 J_1^T)^{-1}$.
    \item Calculate the Null-Space projector of the primary task: $P_1 = (I - J_1^\# J_1)$.
    \item The joint velocity for the primary task is: $\dot{q}_1 = J_1^\# \dot{x}_1$.
    \item Calculate the residual error for the second task: $e_2 = \dot{x}_2 - J_2 \dot{q}_1$.
    \item Project $J_2$ into the null space of $J_1$: $\tilde{J}_2 = J_2 P_1$.
    \item The total solution is: $\dot{q}_{TP} = \dot{q}_1 + \tilde{J}_2^\# e_2$.
\end{enumerate}


\begin{examtip}
If the augmented Jacobian $J_A = \begin{bmatrix} J_1 \\ J_2 \end{bmatrix}$ is FULL RANK, then the Task Priority solution exactly matches the Task Augmentation (TA) solution: $\dot{q}_{TA} = J_A^\# \begin{bmatrix} \dot{x}_1 \\ \dot{x}_2 \end{bmatrix}$. Priority only matters when tasks conflict (i.e., algorithmic singularity).
\end{examtip}

% ---------------------------------------------------------
\section{Type A.2: Projected and Reduced Gradient}

{Algorithm: Projected Gradient (PG)}
\textbf{Goal:} Execute a task $\dot{x}$ while maximizing/minimizing a scalar objective function $H(q)$ (e.g., distance from joint limits).
\begin{enumerate}
    \item Compute the unconstrained solution: $\dot{q}_{task} = J^\# \dot{x}$.
    \item Compute the gradient of your objective function: $\nabla_q H = \begin{bmatrix} \frac{\partial H}{\partial q_1} & \frac{\partial H}{\partial q_2} & \dots \end{bmatrix}^T$.
    \item Define desired null-space velocity: $\dot{q}_0 = \pm \alpha \nabla_q H$ (use $+$ to maximize, $-$ to minimize, with $\alpha > 0$).
    \item Calculate Null-Space projector: $P = (I - J^\# J)$.
    \item Total velocity: $\dot{q}_{PG} = \dot{q}_{task} + P \dot{q}_0$.
\end{enumerate}


{Algorithm: Reduced Gradient (RG)}
\textbf{Goal:} Same as PG, but computationally lighter (no pseudoinverse).
\begin{enumerate}
    \item Partition the Jacobian into a square non-singular matrix $J_a$ ($m \times m$) and $J_b$ ($m \times (n-m)$): $J = [J_a \mid J_b]$.
    \item Partition joints accordingly: $q_a$ (dependent) and $q_b$ (independent).
    \item Compute the \textit{Reduced Gradient}: $\nabla_{q_b} H' = [-(J_a^{-1} J_b)^T \mid I] \nabla_q H$.
    \item Update independent joints: $\dot{q}_b = \pm \alpha \nabla_{q_b} H'$.
    \item Solve for dependent joints to ensure task execution: $\dot{q}_a = J_a^{-1}(\dot{x} - J_b \dot{q}_b)$.
\end{enumerate}

% ---------------------------------------------------------
\section{Type B: SNS Method (Hard Bounds)}

{Algorithm: Saturation in the Null Space (SNS)}
\textbf{Goal:} Find joint velocities that realize the task but strictly respect bounds $|\dot{q}_i| \leq V_{max}$.
\begin{enumerate}
    \item \textbf{Initial guess:} Calculate unconstrained solution $\dot{q} = J^\# \dot{x}$.
    \item \textbf{Check bounds:} Does any element of $\dot{q}$ exceed its maximum? If NO, done. If YES, proceed to 3.
    \item \textbf{Saturate:} Take the joint $k$ that most violates the bound. Set $\dot{q}_k = V_{max} \times \text{sign}(\dot{q}_k)$.
    \item \textbf{Reduce system:} Remove column $k$ from $J$ to get $J_{-k}$. Modify the task: $\dot{x}_{new} = \dot{x} - J_k \dot{q}_k$.
    \item \textbf{Re-solve:} Calculate new velocity for remaining joints: $\dot{q}_{-k} = J_{-k}^\# \dot{x}_{new}$.
    \item \textbf{Re-check:} Check if the new combined $\dot{q}$ satisfies all bounds. If not, repeat.
\end{enumerate}

% ---------------------------------------------------------
\section{Type C: Deriving the Dynamic Model}

{Algorithm: Deriving the Inertia Matrix M}
\textbf{Goal:} Find the symmetric matrix $M(q)$ from kinetic energy $T = \frac{1}{2}\dot{q}^T M(q) \dot{q}$.
\begin{enumerate}
    \item For each link $i$, compute the linear velocity of its Center of Mass, $v_{ci}$, and its angular velocity, $\omega_i$.
    \item Write Kinetic Energy for each link: $T_i = \frac{1}{2}m_i \|v_{ci}\|^2 + \frac{1}{2}\omega_i^T I_{ci} \omega_i$.
    \item Sum them up: $T_{tot} = \sum T_i$.
    \item Expand $T_{tot}$ and group terms into the quadratic form $\frac{1}{2} \sum \sum m_{ij}(q)\dot{q}_i \dot{q}_j$. The coefficients form $M(q)$.
\end{enumerate}

{Algorithm: Deriving Coriolis and Centrifugal terms c}
\textbf{Method: Christoffel Symbols (Fast Matrix form)}
Let $M_k$ be the $k$-th column of $M(q)$. The $k$-th element of vector $c$ is $c_k = \dot{q}^T C_k(q) \dot{q}$, where:
\begin{equation}
    C_k(q) = \frac{1}{2} \left[ \frac{\partial M_k}{\partial q} + \left(\frac{\partial M_k}{\partial q}\right)^T - \frac{\partial M}{\partial q_k} \right]
\end{equation}
Stack the resulting scalars to get the vector $c(q, \dot{q})$.

{Algorithm: Deriving Gravity Vector g}
\begin{enumerate}
    \item Find the height $y_{ci}$ (component aligned with gravity) of each CoM.
    \item Write total potential energy: $U(q) = \sum m_i g_0 y_{ci}$.
    \item Differentiate: $g(q) = \left( \frac{\partial U(q)}{\partial q} \right)^T$.
\end{enumerate}

% ---------------------------------------------------------
\section{Type D: Skew-Symmetry and Factorizations}

{Algorithm: Finding Factorizations S1 and S2}
\textbf{Goal:} Find matrices $S$ such that $S \dot{q} = c$, one making $\dot{M} - 2S$ skew-symmetric ($S_1$), and one not ($S_2$).
\begin{enumerate}
    \item \textbf{Skew-symmetric $S_1$:} Stack the Christoffel matrices you found earlier!
    \begin{equation}
        S_1(q,\dot{q}) = \begin{bmatrix} \dot{q}^T C_1(q) \\ \dot{q}^T C_2(q) \\ \dots \end{bmatrix}
    \end{equation}
    
    \item \textbf{NON skew-symmetric $S_2$:} 
    Add a dummy matrix $S_0(\dot{q})$ to $S_1$ such that $S_0(\dot{q})\dot{q} = 0$. For a 2-DOF robot, choose:
    $$ S_0(\dot{q}) = \begin{bmatrix} 0 & -\dot{q}_2 \\ \dot{q}_2 & 0 \end{bmatrix} $$
    Then $S_2 = S_1 + S_0$. Since $S_0 \dot{q} = 0$, $S_2 \dot{q} = c(q, \dot{q})$, but $\dot{M}-2S_2$ will NOT be skew-symmetric.
\end{enumerate}

% ---------------------------------------------------------
\section{Type E: Linear Parameterization and Identification}

{Algorithm: Extracting Regressor Y and parameters a}
\textbf{Goal:} Rewrite $M\ddot{q} + c + g = \tau$ as $Y(q, \dot{q}, \ddot{q})\mathbf{a} = \tau$.
\begin{enumerate}
    \item Identify all blocks of constants (masses, lengths, inertias). Group terms that always appear together (e.g., $I_1 + m_1 l_{c1}^2$) into a single dynamic coefficient $a_k$.
    \item Define your vector $\mathbf{a} = [a_1, a_2, \dots, a_p]^T$.
    \item Rewrite the equations of motion isolating the variables.
    \item Extract the variables into the $n \times p$ matrix $Y$. 
\end{enumerate}

{Algorithm: Dynamic Identification}
Once you have $Y\mathbf{a} = \tau$, to find the true parameters $\mathbf{a}$ experimentally:
\begin{enumerate}
    \item Sample $q, \dot{q}, \ddot{q}$ and $\tau$ at $N$ time instants along an exciting trajectory.
    \item Build the stacked Regressor Matrix $\bar{Y}$ ($nN \times p$) and Torque vector $\bar{\tau}$ ($nN \times 1$).
    \item Solve using Least Squares: $\hat{\mathbf{a}} = \bar{Y}^\# \bar{\tau} = (\bar{Y}^T \bar{Y})^{-1} \bar{Y}^T \bar{\tau}$.
\end{enumerate}

% ---------------------------------------------------------
\section{Type F: Dynamic Redundancy Resolution}

{Algorithm: Minimizing Torque Norm}
\textbf{Goal:} Find the joint acceleration $\ddot{q}$ (or torque $\tau$) that executes a Cartesian task $\ddot{x}$ while minimizing the norm of the joint torques $\|\tau\|^2$.
\begin{enumerate}
    \item Write the task at acceleration level: $J\ddot{q} = \ddot{x} - \dot{J}\dot{q}$.
    \item Substitute the dynamic model $\ddot{q} = M^{-1}(\tau - c - g)$ into the task equation.
    \item The problem becomes finding $\tau$ to satisfy the task. The minimum norm solution is obtained using the inertia-weighted pseudoinverse:
    \begin{equation}
        \tau_{min} = (J M^{-1})^\# (\ddot{x} - \dot{J}\dot{q} + J M^{-1}(c + g))
    \end{equation}
\end{enumerate}

% ---------------------------------------------------------
\section{Type G: Time Scaling under Torque Bounds}

{Algorithm: Uniform Time Scaling}
\textbf{Goal:} Find scale factor $k \ge 1$ to satisfy torque bounds $|u_i| \le U_{max}$ without altering the geometric path.
\begin{enumerate}
    \item Determine the max required torque $\tau_{orig}(t)$ during the given trajectory using Inverse Dynamics.
    \item Separate the torque into Gravity terms ($g$) which DO NOT scale, and Inertial/Coriolis terms ($\tau_{dyn}$) which scale with $1/k^2$.
    \item Set up the inequality for the most critical joint:
    \begin{equation}
        \left| \frac{\tau_{dyn}}{k^2} + g \right| \leq U_{max}
    \end{equation}
    \item Solve for $k$:
    \begin{equation}
        k \geq \sqrt{\frac{|\tau_{dyn}|}{U_{max} - |g|}}
    \end{equation}
    \item The new feasible execution time is $T_s = k \cdot T_{old}$.
\end{enumerate}

\end{document}